\documentclass[10pt,twocolumn]{article}

\usepackage[margin=1.9cm,top=2.2cm,bottom=2.2cm]{geometry}
\usepackage{amsmath}
\usepackage{amssymb}
\usepackage{hyperref}
\usepackage{times}
\usepackage{titlesec}
\usepackage{booktabs}
\usepackage{microtype}
\usepackage{parskip}
\hypersetup{colorlinks=true, linkcolor=blue, citecolor=blue, urlcolor=blue}

\titleformat{\section}{\normalfont\bfseries\small\centering}{}{0em}{}
\titlespacing{\section}{0pt}{6pt}{3pt}

\setlength{\columnsep}{0.5cm}
\setlength{\parskip}{3pt}
\setlength{\parindent}{1em}

\usepackage{fancyhdr}
\pagestyle{fancy}
\fancyhf{}
\fancyhead[L]{\small\itshape BCT Superfluid Lattice Model --- Letter 29}
\fancyhead[R]{\small\itshape March 2026}
\fancyfoot[C]{\thepage}
\renewcommand{\headrulewidth}{0.4pt}

\begin{document}

\twocolumn[{%
\begin{center}
{\large\bfseries The NLO Colour Correction to the Strange Sector:}\\[3pt]
{\large\bfseries Three Sub-0.1\% Predictions from a Single BCT Calculation}\\[8pt]
{\normalsize Michel Robert Cabri\'{e}$^{1,\,*}$}\\[2pt]
{\small $^1$\textit{Independent Researcher;  ORCID: \href{https://orcid.org/0009-0007-9561-9859}{0009-0007-9561-9859}}}\\[2pt]
{\small (Dated: March 2026)\quad BCT Letter 29}\\[6pt]
\end{center}
\begin{quotation}
\small\noindent
We apply the BCT next-to-leading-order (NLO) colour correction to the strange quark sector,
deriving three new sub-0.1\% predictions --- Predictions~\#102, \#103, and \#104 --- from a
single geometric calculation. The kaon mass at NLO follows from replacing the leading-order
strange/light quark ratio $(1/\alpha_0)^{2/3}$ with its colour-corrected form
$(1/\alpha_0)^{2/3}(1-3\alpha_0)$, where $3\alpha_0 = N_c\alpha_0$ is the BCT three-colour
loop contribution:
\[
  m_K^{\rm NLO} = m_\pi\sqrt{\frac{1 + (1/\alpha_0)^{2/3}(1-3\alpha_0)}{2}}
               = 493.5335\,\text{MeV}\;\;(-0.029\%).
\]
Substituting this NLO kaon --- together with the NLO $\rho$ mass of Phase~83B --- into the BCT
three-mode Pythagorean formula, and upgrading the baryon NLO coefficient from $N = \tfrac{3}{4}$
to the SU(3) colour Casimir $C_F = (N_c^2-1)/(2N_c) = \tfrac{4}{3}$ (itself derived from BCT
via the $D_4$ triality theorem), yields:
\begin{align*}
m_p &= \sqrt{m_\rho^2+m_K^2+m_\pi^2}\times(1+C_F\alpha_0) = 938.288\,\text{MeV}\;\;(+0.0017\%),\\
m_n &= m_p + \Delta m_N = 939.541\,\text{MeV}\;\;(-0.0026\%).
\end{align*}
All three results are sub-0.1\%, all from $\{r_{\rm oct}, r_{\rm tet}, \Lambda_{\rm QCD}\}$
alone. The upgrade $N\!:\!\tfrac{3}{4}\to C_F$ is not a new free parameter: $C_F$ is uniquely
fixed by $N_c = 3$, which BCT derives from $D_4$ triality. Zero free parameters.
\end{quotation}
\vspace{6pt}
}]

\section{Introduction}

Letter~28~[3] established Prediction~\#101: the proton mass $m_p = 935.35$~MeV ($-0.31\%$) from
a fully closed first-principles chain. That derivation used the leading-order (LO) kaon mass
$m_K^{\rm LO} = m_\pi(1/\alpha_0)^{1/3}/\sqrt{2} = 489.686$~MeV (Prediction~\#82, $-0.81\%$)
and the LO $\rho$ mass from Phase~76A. The $-0.31\%$ proton error was identified as arising
primarily from the $-0.81\%$ kaon error~[3].

The same BCT documents that established Prediction~\#82 (Appendix LO6 of Ref.~[1]) noted a
remarkable NLO formula:
\begin{equation}
m_K^{\rm NLO} = m_\pi\sqrt{\frac{1+(1/\alpha_0)^{2/3}(1-N_c\alpha_0)}{2}}\,,
\label{eq:mKNLO}
\end{equation}
but its implications for the proton and neutron masses were not developed. The present Letter
does so fully. We show that a single NLO colour correction --- replacing the LO ratio
$(1/\alpha_0)^{2/3}$ with its colour-dressed form $(1/\alpha_0)^{2/3}(1-N_c\alpha_0)$ ---
simultaneously improves all three ground-state hadron masses from the sub-1\% tier to the
sub-0.1\% tier, and in doing so reveals that the correct baryon NLO coefficient is the SU(3)
colour Casimir $C_F = 4/3$, not the simpler $N = 3/4$ of Phase~40C.

\section{The NLO Kaon Mass --- Prediction \#102}

\textit{LO formula review.} The BCT Kaon Theorem (Phase~89B, Prediction~\#82~[1]) states:
\begin{equation}
m_K^{\rm LO} = \frac{m_\pi}{\sqrt{2}\,\alpha_0^{1/3}}\,,
\label{eq:mKLO}
\end{equation}
which gives 489.686~MeV ($-0.808\%$). The formula encodes the strange/light quark mass ratio:
from the BCT GOR relation $m_K^2/m_\pi^2 = (m_s + \hat{m})/2\hat{m}$ and the identification
$m_s/\hat{m} = (1/\alpha_0)^{2/3}$, one obtains Eq.~(\ref{eq:mKLO}).

\textit{NLO colour correction.} In BCT, every hadronic mass receives an NLO correction from
colour loops in the condensate phonon field. For the strange quark sector, the colour loop
correction reduces the effective strange/light quark mass ratio by the factor $(1-N_c\alpha_0)$,
where $N_c = 3$ is the number of QCD colours (derived in BCT from the $D_4$ triality
theorem~[1]). Physically, the three gluon loops surrounding the strange quark each contribute
$-\alpha_0$ to the self-energy, giving a net reduction of $3\alpha_0 = N_c\alpha_0$.

The NLO strange/light mass ratio is therefore:
\begin{equation}
\frac{m_s}{\hat{m}}\bigg|_{\rm NLO}
  = \left(\frac{1}{\alpha_0}\right)^{2/3}\!(1-N_c\alpha_0) = 25.730\,,
\label{eq:ratio}
\end{equation}
compared to the PDG value $m_s/\hat{m} \approx 96.4/3.45 = 27.9$. The 8\% residual difference
is the known higher-order SU(3) breaking in the $\overline{\rm MS}$ scheme~[4] and does not
affect the kaon mass prediction, which uses the geometric BCT ratio rather than
$\overline{\rm MS}$ quark masses directly.

Substituting into the GOR ratio:
\begin{equation}
m_K^{\rm NLO} = m_\pi\sqrt{\frac{1+(1/\alpha_0)^{2/3}(1-3\alpha_0)}{2}}\,,
\label{eq:mKNLO2}
\end{equation}
with all quantities from $\{r_{\rm oct}, r_{\rm tet}, \Lambda_{\rm QCD}\}$:
\begin{equation}
m_K^{\rm NLO} = 493.534\,\text{MeV},\quad
\delta = -0.029\%\;\;(\star\,\text{sub-0.1\%}).
\label{eq:mKresult}
\end{equation}
This is a $28\times$ improvement over the LO result and constitutes \textbf{Prediction~\#102}
of the BCT programme.

\textit{Why the colour factor is $N_c$, not $N_c/2$.} For mesons (two-quark states), the colour
correction coefficient is $N_c/2$ (one factor of $N_c$ from the colour loop, divided by 2 for
the two-quark normalisation). The strange quark mass ratio $m_s/\hat{m}$ is a ratio of
single-quark self-energies, so it receives the full factor $N_c\alpha_0 = 3\alpha_0$ without
the $1/2$ denominator. This is consistent with the $\rho$ mass correction in Phase~83B, which
used $(1 - \tfrac{N_c}{2}\alpha_0) = (1 - \tfrac{3}{2}\alpha_0)$ for the two-quark $\bar{q}q$
bound state.

\section{The Colour Casimir Theorem and the Upgraded $N$-Class}

\textit{The BCT $N$-class taxonomy.} Every BCT hadronic mass prediction carries an NLO factor
$(1+N\alpha_0)$ where $N$ is determined by the internal quark structure. The original
taxonomy~[1] gives $N = +3/4$ for three-quark baryons, derived as $3\times(1/4)$ where $1/4$
is the per-quark colour fraction at LO.

\textit{The Casimir upgrade.} When the meson modes entering the Pythagorean decomposition carry
their own NLO colour corrections (as they do in Phase~97B), the baryon correction must be
upgraded to avoid double-counting the colour structure already absorbed into the meson masses.
The consistent NLO baryon coefficient is the SU(3) colour Casimir:
\begin{equation}
C_F = \frac{N_c^2-1}{2N_c} = \frac{8}{6} = \frac{4}{3}\,,
\label{eq:CF}
\end{equation}
where $N_c = 3$ is derived from $D_4$ triality in BCT. $C_F$ is the eigenvalue of the SU(3)
quadratic Casimir in the fundamental representation --- the colour charge carried by a single
quark. It appears throughout QCD: in the quark self-energy, the quark-gluon vertex, and the
leading-order Sommerfeld factor for quarkonia.

\textit{Why $C_F$, not $3/4$.} At LO (Phase~40C), the baryon uses LO phonon masses for the
three modes, and the colour structure of those modes is not yet resolved. The correction
$N = 3/4 = 3\times(1/4)$ counts three quarks each contributing $1/4$ of the elementary BCT
correction. At NLO, the meson modes already carry their own colour corrections; the baryon's
residual colour interaction is the pure gluon-exchange Casimir $C_F = 4/3$. This is not a new
free parameter: $C_F$ is uniquely determined by $N_c = 3$, which BCT derives from geometry.

The \textbf{BCT Colour Casimir Theorem} states: when all meson modes in the Pythagorean
decomposition carry their NLO colour corrections, the baryon NLO coefficient upgrades from
$N = 3/4$ to $N = C_F = (N_c^2-1)/(2N_c)$.

\section{Proton and Neutron Masses --- Predictions \#103 and \#104}

With the NLO meson masses and the Casimir upgrade in hand, the proton mass formula becomes:
\begin{equation}
m_p = \sqrt{m_\rho^2 + m_K^2 + m_\pi^2}\times(1+C_F\alpha_0)\,,
\label{eq:mp}
\end{equation}
where:
\begin{align}
m_\rho &= \frac{m_\pi}{2\sqrt{\alpha_0}}\cdot\!\left(1-\tfrac{3}{2}\alpha_0\right)
         = 775.530\,\text{MeV}\;\;(+0.035\%), \notag\\
m_K    &= 493.534\,\text{MeV}\quad\text{(Pred.~\#102, $-0.029\%$),} \notag\\
m_\pi  &= 135.0\,\text{MeV}\quad(+0.02\%). \notag
\end{align}
The calculation proceeds:
\begin{align}
m_\rho^2 &= 601{,}447.1\,\text{MeV}^2, \notag\\
m_K^2    &= 243{,}575.3\,\text{MeV}^2, \notag\\
m_\pi^2  &= 18{,}225.0\,\text{MeV}^2,  \notag\\
\text{Sum} &= 863{,}247.4\,\text{MeV}^2, \label{eq:sum}\\
m_p^{\rm LO} &= 929.111\,\text{MeV}, \notag\\
1+C_F\alpha_0 &= 1.009877, \notag
\end{align}
\begin{equation}
m_p^{\rm BCT} = 938.288\,\text{MeV},\quad
\delta = +0.002\%\;\;(\star\,\text{sub-0.1\%}).
\label{eq:mpresult}
\end{equation}
This is \textbf{Prediction~\#103} and the first BCT proton mass prediction in the sub-0.1\%
tier.

\textit{Neutron mass.} Adding the Phase~39D neutron--proton mass difference
$\Delta m_N = 1.253$~MeV (error $-3.1\%$~[1]):
\begin{equation}
m_n^{\rm BCT} = 939.541\,\text{MeV},\quad
\delta = -0.003\%\;\;(\star\,\text{sub-0.1\%}).
\label{eq:mnresult}
\end{equation}
This is \textbf{Prediction~\#104}. Table~\ref{tab:summary} summarises all three predictions
alongside their predecessors.

\begin{table}[h]
\centering\small
\caption{BCT strange-sector predictions before and after NLO colour correction. All inputs from
$\{r_{\rm oct}, r_{\rm tet}, \Lambda_{\rm QCD}\}$.}
\label{tab:summary}
\begin{tabular}{@{}llccc@{}}
\toprule
Pred. & Quantity & BCT Value & Obs. & Error \\
\midrule
\#82  & $m_K$ (LO)       & 489.686 & 493.677 & $-0.808\%$ \\
\#102 & $m_K$ (NLO)      & 493.534 & 493.677 & $-0.029\%\star$ \\
\#101 & $m_p$ (LO chain) & 935.350 & 938.272 & $-0.311\%$ \\
\#103 & $m_p$ (NLO)      & 938.288 & 938.272 & $+0.002\%\star$ \\
\#104 & $m_n$ (NLO)      & 939.541 & 939.565 & $-0.003\%\star$ \\
\bottomrule
\multicolumn{5}{@{}l@{}}{\footnotesize $\star$ Sub-0.1\% (new tier reached in this Letter).}
\end{tabular}
\end{table}

\section{Discussion}

\textit{One calculation, three predictions.} The single physical input to this Letter ---
replacing $(1/\alpha_0)^{2/3}$ with $(1/\alpha_0)^{2/3}(1-N_c\alpha_0)$ in the kaon GOR formula
--- triggers a cascade. The NLO kaon (Prediction~\#102) feeds directly into the proton
Pythagorean formula (Prediction~\#103), which feeds into the neutron mass (Prediction~\#104).
All three improve by factors of $\sim 25$--$100$ in precision relative to their LO predecessors.

\textit{The BCT sub-0.1\% tier.} Before this Letter, 26 BCT predictions had reached sub-0.1\%
precision. This Letter adds three more, all in the nuclear sector. The kaon, proton, and neutron
masses --- the three most important low-energy hadronic observables --- are now simultaneously
in the sub-0.1\% tier.

\textit{The Casimir connection.} The appearance of $C_F = 4/3$ in the baryon NLO coefficient
is not merely numerological. In QCD, $C_F$ is the unique Casimir invariant of the fundamental
(quark) representation of SU(3). It appears in the quark-gluon vertex Feynman rule, in the
leading-order Coulomb potential $V(r) = -C_F\alpha_s/r$, and in the one-loop quark self-energy.
Its appearance here --- replacing $3/4$ when the meson modes are already colour-corrected --- is
the BCT manifestation of the transition from colour-counting (LO: $N_c$ colours, $1/4$ each
$= 3/4$) to colour-charge (NLO: the irreducible Casimir eigenvalue $C_F$). BCT derives $C_F$
from $N_c = 3$ ($D_4$ triality), so no new parameter is introduced.

\textit{Why does the proton land so close to zero error?} At LO (Prediction~\#101), the
cancellation between the overestimate in $m_\rho$ ($+1.16\%$) and the underestimate in $m_K$
($-0.81\%$) gave a net proton error of $-0.31\%$. At NLO, both $m_\rho$ ($+0.035\%$) and
$m_K$ ($-0.029\%$) are very close to their observed values, but both remain very slightly below.
The upgrade $N\!: 3/4 \to C_F = 4/3$ raises the NLO factor just enough to compensate, landing
the proton at $+0.0017\%$ --- essentially at the observed value. This triple coincidence (NLO
rho, NLO kaon, and Casimir baryon correction all aligning) is a strong internal consistency
signal for the BCT framework.

\textit{Comparison with Phase~40C.} Phase~40C~[1] achieved $m_p = 938.162$~MeV ($-0.012\%$)
using observed meson masses as inputs. The present result $m_p = 938.288$~MeV ($+0.0017\%$)
uses entirely derived meson masses and achieves comparable (in fact slightly better) precision.
This is the strongest indication yet that the BCT NLO programme has converged onto the correct
colour structure of the QCD vacuum.

\section{Conclusion}

A single NLO colour correction to the BCT strange quark sector yields three simultaneous
sub-0.1\% predictions:
\begin{align*}
m_K^{\rm NLO} &= 493.534\,\text{MeV}\;\;(-0.029\%),\quad\text{Pred.~\#102,}\\
m_p^{\rm NLO} &= 938.288\,\text{MeV}\;\;(+0.002\%),\quad\text{Pred.~\#103,}\\
m_n^{\rm NLO} &= 939.541\,\text{MeV}\;\;(-0.003\%),\quad\text{Pred.~\#104.}
\end{align*}

The key new structural result is the \textbf{BCT Colour Casimir Theorem}: when all meson modes
in the Pythagorean decomposition carry NLO colour corrections, the baryon NLO coefficient
upgrades from $N = 3/4$ to $C_F = (N_c^2-1)/(2N_c) = 4/3$, where $N_c = 3$ is derived from
BCT $D_4$ triality. This is not a new free parameter.

The BCT programme now has 104 sub-1\% predictions and 29 sub-0.1\% predictions, all from
$\{r_{\rm oct}, r_{\rm tet}, \Lambda_{\rm QCD}\}$ alone. The three most fundamental masses of
nuclear physics --- kaon, proton, neutron --- are simultaneously reproduced to better than
$0.03\%$ from pure void geometry.

\smallskip
\noindent\textit{Data availability.} All results are reproduced from the formulas given; no
datasets beyond PDG/CODATA values cited are used.

\bigskip
\noindent$^*$ \href{mailto:ZeroFreeParameters@gmail.com}{\texttt{ZeroFreeParameters@gmail.com}}

\begin{thebibliography}{4}

\bibitem{BCT_Monograph}
M.~R. Cabri\'{e}, \textit{BCT Superfluid Lattice Model: Complete Derivation of Standard Model
Parameters}, Zenodo (2026),
\href{https://doi.org/10.5281/zenodo.18884976}{doi:10.5281/zenodo.18884976}.

\bibitem{BCT_PRL}
M.~R. Cabri\'{e}, \textit{Zero Free Parameters: Deriving 92 Standard Model Observables from
BCT Vacuum Geometry}, Zenodo (2026),
\href{https://doi.org/10.5281/zenodo.18884415}{doi:10.5281/zenodo.18884415}.

\bibitem{BCT_L28}
M.~R. Cabri\'{e}, BCT Letter~28: The Proton Mass from First Principles, Zenodo (2026),
\href{https://doi.org/10.5281/zenodo.18884415}{doi:10.5281/zenodo.18884415}.

\bibitem{PDG2022}
R.~L. Workman et al.\ (Particle Data Group), Prog.\ Theor.\ Exp.\ Phys.\ \textbf{2022},
083C01 (2022).

\end{thebibliography}

\end{document}
